\documentclass[11pt]{article}
\usepackage[margin=1in]{geometry}
\usepackage{amsmath,amssymb}
\usepackage[hidelinks]{hyperref}

\title{Literature Status: Monotonicity of \texorpdfstring{$FL^p$}{FLp} and \texorpdfstring{$a$-$FL^p$}{a-FLp}}
\author{FA/Banach run \texttt{fa\_banach\_001}}
\date{2026-07-18}

\begin{document}
\maketitle

\section*{Status}

This is a literature-already-answered packet, not an original proof from the
run.  The first displayed question in Section ``Current developments and open
questions'' of Chatterji--Drutu--Haglund, arXiv:0704.3749, asks whether, for
every $p\ge q\ge 2$,
\[
  FL^p \Longrightarrow FL^q
  \qquad\text{and}\qquad
  a\text{-}FL^q \Longrightarrow a\text{-}FL^p .
\]
Marrakchi--de la Salle, arXiv:2001.02490, answer this question
affirmatively.

\section*{Original Question}

The original source is:
\begin{quote}
I. Chatterji, C. Drutu, and F. Haglund,
\emph{Kazhdan and Haagerup properties from the median viewpoint},
arXiv:0704.3749.
\end{quote}
Definition 1.4 says that a topological group $G$ has property $FL^p$ when
every continuous affine isometric action of $G$ on a space $L^p(X,\mu)$ has
bounded orbits, and that $G$ is $a$-$FL^p$-menable when it has a proper
continuous affine isometric action on some $L^p(X,\mu)$.

The first question in Section ``Current developments and open questions''
asks whether Corollary 3 generalizes to the monotonicity statement above.

\section*{Later Answer}

The supporting source is:
\begin{quote}
A. Marrakchi and M. de la Salle,
\emph{Isometric actions on $L_p$-spaces: dependence on the value of $p$},
arXiv:2001.02490.
\end{quote}
Its abstract explicitly says that it answers a question by
Chatterji--Drutu--Haglund.  Theorem 1 states that if $0<p\le q<\infty$, then
every continuous affine isometric action $\alpha:G\curvearrowright L_p$ gives
a continuous affine isometric action $\beta:G\curvearrowright L_q$ such that
\[
  \|\alpha_g(0)\|_{L_p}^p = \|\beta_g(0)\|_{L_q}^q
  \qquad(g\in G).
\]
Consequently unbounded orbits, and metrically proper orbits, pass from
exponent $p$ to every larger exponent $q$.

Corollary 2 records this in interval form: for every topological group $G$,
the exponents admitting a continuous affine isometric $L_p$ action with
unbounded orbits form $(p_G,\infty)$ or $[p_G,\infty)$, and the exponents
admitting a proper continuous affine isometric $L_p$ action form
$(p'_G,\infty)$ or $[p'_G,\infty)$.

\section*{Translation}

The negation of $FL^r$ is the existence of a continuous affine isometric
$L^r$ action with unbounded orbits.  Since the set of such exponents is an
upper interval, its complement is downward closed.  Hence if $p\ge q\ge 2$
and $G$ has $FL^p$, then $G$ has $FL^q$.

Similarly, $a$-$FL^r$-menability is exactly the existence of a proper
continuous affine isometric action on some $L^r$ space.  Since the proper
action exponents form an upper interval, $a$-$FL^q$-menability implies
$a$-$FL^p$-menability for every $p\ge q\ge 2$.

\section*{Verification Notes}

This is not a same-paper extraction false positive: arXiv:0704.3749 asks the
question in 2007/2009, while arXiv:2001.02490 is a later independent paper.
The supporting paper explicitly identifies the Chatterji--Drutu--Haglund
question and says that Corollary 2 answers it.  Its introduction also states
that throughout the paper an $L_p$ space means $L_p(X,\mu)$ for a standard
measure space, matching the formulation of $FL^p$ in the original source.

\section*{Limitations}

This packet only records the literature answer to the monotonicity question.
It does not compute the constants $p_G$ or $p'_G$, and it does not settle the
separate question in arXiv:0704.3749 about equivalence of different large
$FL^p$ properties or prescribed threshold examples.

\begin{thebibliography}{9}
\bibitem{CDH}
I. Chatterji, C. Drutu, and F. Haglund,
\emph{Kazhdan and Haagerup properties from the median viewpoint},
arXiv:0704.3749.

\bibitem{MdS}
A. Marrakchi and M. de la Salle,
\emph{Isometric actions on $L_p$-spaces: dependence on the value of $p$},
arXiv:2001.02490.
\end{thebibliography}

\end{document}
